% Preámbulo común de las Tareas de Teoría de la Computación (XeLaTeX)
\usepackage{fontspec}
\IfFontExistsTF{Times New Roman}{\setmainfont{Times New Roman}}{\setmainfont{TeX Gyre Termes}}
\setmonofont{DejaVu Sans Mono}[Scale=0.85]
\newfontfamily\simbolos{DejaVu Sans}[Scale=0.9]
\usepackage[spanish,mexico,es-noshorthands]{babel}
\usepackage[letterpaper,top=2.6cm,bottom=2.2cm,left=2.2cm,right=2.2cm,headheight=28pt]{geometry}
\usepackage{amsmath,amssymb,amsthm}
\usepackage{graphicx}
\usepackage[table]{xcolor}
\usepackage{booktabs,longtable,array,multicol,tabularx}
\usepackage{enumitem}
\setlist{nosep,leftmargin=*}
\usepackage{fancyhdr}
\usepackage[most]{tcolorbox}
\usepackage[hidelinks]{hyperref}
\usepackage{newunicodechar}
\newunicodechar{λ}{\ensuremath{\lambda}}
\newunicodechar{Σ}{\ensuremath{\Sigma}}
\newunicodechar{δ}{\ensuremath{\delta}}
\newunicodechar{∅}{\ensuremath{\varnothing}}
\newunicodechar{→}{\ensuremath{\to}}
\newunicodechar{…}{\ldots}
\newunicodechar{·}{\ensuremath{\cdot}}
\newunicodechar{≡}{\ensuremath{\equiv}}
\newunicodechar{≥}{\ensuremath{\ge}}
\newunicodechar{≤}{\ensuremath{\le}}
\newunicodechar{∈}{\ensuremath{\in}}
\newunicodechar{ⁿ}{\textsuperscript{\textit{n}}}
\newunicodechar{ᵐ}{\textsuperscript{\textit{m}}}
\newunicodechar{⁺}{\textsuperscript{+}}
\newunicodechar{₀}{\textsubscript{0}}
\newunicodechar{₁}{\textsubscript{1}}
\newunicodechar{₂}{\textsubscript{2}}
\newunicodechar{₃}{\textsubscript{3}}
\newunicodechar{✓}{{\simbolos ✓}}
\newunicodechar{✗}{{\simbolos ✗}}
\newunicodechar{⁻}{\textsuperscript{−}}
\newunicodechar{¹}{\textsuperscript{1}}
\newunicodechar{²}{\textsuperscript{2}}
\newunicodechar{³}{\textsuperscript{3}}

\newunicodechar{⁰}{\textsuperscript{0}}
\newunicodechar{⁴}{\textsuperscript{4}}
\newunicodechar{⁵}{\textsuperscript{5}}
\newunicodechar{⁶}{\textsuperscript{6}}
\newunicodechar{⁷}{\textsuperscript{7}}
\newunicodechar{⁸}{\textsuperscript{8}}
\newunicodechar{⁹}{\textsuperscript{9}}
\newunicodechar{ᵏ}{\textsuperscript{\textit{k}}}
\newunicodechar{ᴿ}{\textsuperscript{\textit{R}}}
\newunicodechar{∉}{\ensuremath{\notin}}
\newunicodechar{∣}{\ensuremath{\mid}}
\newunicodechar{∩}{\ensuremath{\cap}}
\newunicodechar{∪}{\ensuremath{\cup}}
\newunicodechar{⊆}{\ensuremath{\subseteq}}
\newunicodechar{⋃}{\ensuremath{\bigcup}}
\newunicodechar{⋯}{\ensuremath{\cdots}}
\newunicodechar{≠}{\ensuremath{\neq}}
\newunicodechar{↦}{\ensuremath{\mapsto}}
\newunicodechar{×}{\ensuremath{\times}}
\newunicodechar{𝒫}{\ensuremath{\mathcal{P}}}
\newunicodechar{↔}{\ensuremath{\leftrightarrow}}
\newunicodechar{⇔}{\ensuremath{\Leftrightarrow}}
\newunicodechar{∋}{\ensuremath{\ni}}
\newunicodechar{⊥}{\ensuremath{\bot}}
\newunicodechar{ᵖ}{\textsuperscript{\textit{p}}}
\newunicodechar{ᵢ}{\textsubscript{\textit{i}}}
\newunicodechar{ⱼ}{\textsubscript{\textit{j}}}
\newunicodechar{ₓ}{\textsubscript{\textit{x}}}
\newunicodechar{ᵃ}{\textsuperscript{\textit{a}}}
\newunicodechar{ᵇ}{\textsuperscript{\textit{b}}}
\newunicodechar{ᶜ}{\textsuperscript{\textit{c}}}
\newunicodechar{∃}{\ensuremath{\exists}}
\newunicodechar{∀}{\ensuremath{\forall}}
\newunicodechar{ˣ}{\textsuperscript{\textit{x}}}
\newunicodechar{ⁱ}{\textsuperscript{\textit{i}}}
\newunicodechar{ₙ}{\textsubscript{\textit{n}}}
\newunicodechar{⟶}{\ensuremath{\longrightarrow}}
\newunicodechar{⇒}{\ensuremath{\Rightarrow}}
\newunicodechar{ʳ}{\textsuperscript{\textit{r}}}
\newunicodechar{ᵗ}{\textsuperscript{\textit{t}}}
\newunicodechar{ʲ}{\textsuperscript{\textit{j}}}
\newunicodechar{ᑫ}{\textsuperscript{\textit{q}}}
\newunicodechar{ᵒ}{\textsuperscript{o}}
\newunicodechar{₄}{\textsubscript{4}}
\newunicodechar{ₖ}{\textsubscript{\textit{k}}}
\newunicodechar{ₘ}{\textsubscript{\textit{m}}}
\newunicodechar{ₚ}{\textsubscript{\textit{p}}}
\newunicodechar{∖}{\ensuremath{\setminus}}
\newunicodechar{¬}{\ensuremath{\neg}}
\newunicodechar{∧}{\ensuremath{\wedge}}
\newunicodechar{∨}{\ensuremath{\vee}}
\newunicodechar{ℕ}{\ensuremath{\mathbb{N}}}
\newunicodechar{Ø}{\ensuremath{\varnothing}}
\newunicodechar{ε}{\ensuremath{\varepsilon}}
\newunicodechar{ω}{\ensuremath{\omega}}
\newunicodechar{π}{\ensuremath{\pi}}
\newunicodechar{τ}{\ensuremath{\tau}}
\newunicodechar{ρ}{\ensuremath{\rho}}
\newunicodechar{Δ}{\ensuremath{\Delta}}
\newunicodechar{⇐}{\ensuremath{\Leftarrow}}
\newunicodechar{−}{\ensuremath{-}}
\newunicodechar{∞}{\ensuremath{\infty}}

\definecolor{guinda}{HTML}{6C1D45}
\definecolor{dorado}{HTML}{B38E5D}
\definecolor{fondo}{HTML}{F7F2F4}

% --- Encabezado en todas las hojas: nombre, boleta y grupo -----------------
\pagestyle{fancy}
\fancyhf{}
\fancyhead[L]{\small\textbf{Rico Espluga Juan Diego}\\\small Boleta: 2025630474 \quad Grupo: 4CV4}
\fancyhead[R]{\small\textcolor{guinda}{Teoría de la Computación}\\\small\NombreTarea}
\fancyfoot[C]{\small\thepage}
\renewcommand{\headrulewidth}{0.6pt}
\fancypagestyle{plain}{}

\setlength{\parindent}{0pt}
\setlength{\parskip}{4pt}
\renewcommand{\arraystretch}{1.15}

% --- Ejercicios ---------------------------------------------------------
% \begin{ejercicio}{97}{*} enunciado \end{ejercicio}
\newtcolorbox{ejercicio}[2]{enhanced,breakable,colback=fondo,colframe=guinda,boxrule=0.6pt,arc=2pt,
  left=6pt,right=6pt,top=4pt,bottom=4pt,fonttitle=\bfseries\small,
  title={Ejercicio #1\ifx&#2&\else\ (obligatorio)\fi}, before skip=10pt, after skip=4pt}
\newcommand{\proc}{\par\textbf{\textcolor{guinda}{Procedimiento.}}\ }
\newcommand{\resp}{\par\textbf{\textcolor{guinda}{Respuesta.}}\ }
\newcommand{\seccion}[1]{\section*{\textcolor{guinda}{#1}}\addcontentsline{toc}{section}{#1}}

% \jflap{archivo sin extensión}{leyenda}: captura de JFLAP o recuadro para insertarla.
\newcommand{\jflap}[2]{%
  \par\begin{center}
  \IfFileExists{img/#1.png}{%
    \includegraphics[width=0.96\linewidth,height=0.4\textheight,keepaspectratio]{img/#1.png}%
  }{%
    \fcolorbox{dorado}{white}{\parbox[c][2.6cm][c]{0.78\linewidth}{\centering
      \textbf{\textcolor{dorado}{Espacio para la captura de JFLAP}}\\[3pt]
      \small Archivo: \texttt{jff/\detokenize{#1}.jff}\\
      \footnotesize Guarde la captura como \texttt{img/\detokenize{#1}.png} y vuelva a compilar.}}%
  }\\[2pt]
  {\small\textit{#2}}
  \end{center}}
